\section{Task-Level Criterion-Association Results}
\label{sec:validity-results}

\subsection{Development freeze}

The first development specification failed before any sealed reveal:
leave-one-world-out improvements were \(0.000421\) for the terminal Brier score
and \(-0.000221\) for the integrated score.  The effect floor was not lowered.
The revised specification changed the exogenous task utility from mode equality
to structural predicates, used task-varying hazards and scalar
selected-versus-alternative contrasts, doubled units per world from 16 to 32,
and retained the stronger layer-aware sensitivity.

Across 24 revised development worlds, the primary event rate was \(0.60384\).
Leave-one-world-out terminal and integrated improvements were \(0.003124\) and
\(0.003927\).  An analytic floor and a 20,000-draw conjunctive Holm power
simulation selected 69 sealed worlds; simulated conjunctive power was
\(0.9208\), with Monte Carlo 95\% lower bound \(0.91698\).

\subsection{Sealed primary analysis}

All \(1{,}242\) sealed fits completed.  The primary population contained
\(26{,}496=69\times32\times3\times4\) unit records (worlds, units, nested
optimizer seeds, and primary arms), with event rate \(0.61632\); both exact controls
retained event rate zero.  Table~\ref{tab:validity-primary} gives the
co-primary analysis.

\begin{table}[ht]
\centering
\caption{Frozen co-primary criterion-validity results across 69 worlds.
Bootstrap and hierarchical columns are simultaneous lower bounds against
zero.}
\label{tab:validity-primary}
\begin{tabular}{lrrrrr}
\toprule
Effect & Mean & World SD & Holm \(p\) & Bootstrap LB & Hierarchical LB \\
\midrule
Terminal Brier
& 0.003519 & 0.009015 & 0.001230 & 0.001428 & 0.001392 \\
Integrated first failure
& 0.003034 & 0.007471 & 0.001230 & 0.001296 & 0.001271 \\
\bottomrule
\end{tabular}
\end{table}

Both means exceed \(\SESOI=0.0025\).  Mean terminal Brier score decreased from
\(0.183661\) to \(0.180143\), and integrated Brier score decreased from
\(0.148317\) to \(0.145283\).  Outcome noncircularity, panel completeness,
event balance, no-refitting, exact-control structural zero, Holm,
cluster-bootstrap, and hierarchical requirements all passed.

The pooled estimand averages four heterogeneous arms.  Table~\ref{tab:validity-arms}
therefore reports the prespecified secondary arm decomposition before any
routing interpretation.

\begin{table}[ht]
\centering
\caption{Arm-level Brier improvements over the frozen scalar model.  These are
secondary descriptive effects, not separately multiplicity-gated estimands.}
\label{tab:validity-arms}
\begin{tabular}{lrr}
\toprule
Arm & Terminal & Integrated \\
\midrule
Layer-routed dense & 0.005402 & 0.003222 \\
Strict factorized & 0.003114 & 0.003077 \\
Random-routed matched sparsity & 0.002445 & 0.002758 \\
Permuted or wrong-routed & 0.003114 & 0.003077 \\
\bottomrule
\end{tabular}
\end{table}

The wrong-routed arm exceeds \(\SESOI\) on both outcomes, and the random-routed
arm exceeds it on the integrated outcome.  Consequently, the pooled criterion
analysis identifies an association between structural augmentation and the
same-task outcomes, not the effect of routing correctness.  The larger
descriptive effects for layer-routed dense than for strict factorization also
do not isolate factorization as the active component.

\subsection{Sensitivity and temporal boundaries}

When direct raw layer mismatches were added to the generic baseline, structural
augmentation improved terminal and integrated Brier scores by \(0.002108\) and
\(0.001467\).  These effects are positive but below \(\SESOI\).  Thus the
primary result establishes an increment beyond a scalar generic baseline, not
a practically nontrivial advantage beyond the prespecified layer-aware
baseline.

Two stronger temporal interpretations were tested only on development data and
failed.  Using the independent probe to predict later task failures yielded
\(-0.000539\) and \(-0.000212\).  Using the first four tasks as calibration to
predict the later four yielded \(-0.002021\) and \(-0.001069\).  Consequently,
the sealed positive result is a same-task criterion association using oracle
candidate endpoints.  It is not evidence for future-task prognosis or an
oracle-free deployment score.
